\documentclass{fiche}
\metadonnees{4}{Le passé du passé}{Bloc 1 — Le récit : présent et passé}

\begin{document}
\entete

\objectifs{%
\textbf{Langue} : pluperfect ; present perfect simple et en \emph{be + -ing} ; lexique du
paysage et de la peur.\par
\textbf{Texte} : Charles Dickens, \emph{Great Expectations} (1861), chapitre \textsc{i}.\par
\textbf{Durée} : 1\,h\,30 — exercices 35 min, texte et questions 30 min, version 25 min.}

%% ===================================================================
\exercice[13 min]{Le pluperfect}

\rappel{\textbf{Rappel.} \emph{Had} + participe passé situe un fait avant un autre fait
passé. Le prétérit raconte, le pluperfect remonte en arrière.}

\consigne{Mettez le verbe entre parenthèses au pluperfect.}

\begin{anglais}
\begin{enumerate}[label=\arabic*.,itemsep=2.2mm,leftmargin=*]
  \item Pip never saw his parents, who \emph{(die)} \trou{had died} long before.
  \item As soon as the man \emph{(empty)} \trou{had emptied} his pockets, he sat him on a
    tombstone.
  \item The convict was covered with mud because he \emph{(walk)} \trou{had walked} through
    the marshes all night.
  \item Pip realised that he \emph{(stay)} \trou{had stayed} in the churchyard too late.
  \item He \emph{(never feel)} \trou{had never felt} so keen a wind in his life.
  \item By the time he reached home, the man \emph{(disappear)} \trou{had disappeared}.
  \item My sister told me that she \emph{(see)} \trou{had seen} a file in the forge.
  \item The iron on his leg showed that he \emph{(escape)} \trou{had escaped} from the
    prison ships.
  \item The clumsy boy \emph{(hardly / sit)} \trou{had hardly sat} down when the man seized
    him again.
  \item After he \emph{(point)} \trou{had pointed} to the village, the man turned him upside
    down.
  \item He was out of breath: he \emph{(run)} \trou{had been running} all the way.
  \item When I looked over my shoulder, he \emph{(already / get)} \trou{had already got} over
    the wall.
\end{enumerate}
\end{anglais}

\note{Phrase 11 : \emph{had been running} plutôt que \emph{had run}, parce que c'est
l'activité et sa durée qui expliquent l'essoufflement, non son résultat. Même partage qu'à
l'exercice 2, mais reporté d'un cran dans le passé.}

\note[Deux tours à retenir]{\emph{As soon as} et \emph{after} (2, 10) appellent
naturellement le pluperfect : le premier fait est achevé avant que le second commence.
\emph{Hardly\ldots when} (9) et \emph{no sooner\ldots than} servent à dire qu'une action en
a immédiatement suivi une autre ; la langue soutenue inverse alors sujet et auxiliaire :
\emph{Hardly had he sat down when\ldots}}

%% ===================================================================
\exercice[10 min]{Present perfect simple ou en \emph{be + -ing}}

\rappel{\textbf{Rappel.} Le present perfect en \emph{be + -ing} porte sur l'activité et sa
durée ; le present perfect simple sur le résultat acquis ou le compte.}

\consigne{Mettez le verbe entre parenthèses à la forme qui convient.}

\begin{anglais}
\begin{enumerate}[label=\arabic*.,itemsep=2.2mm,leftmargin=*]
  \item I \emph{(read)} \trou{have read} this book three times.
  \item I \emph{(read)} \trou{have been reading} for two hours; my eyes hurt.
  \item Look at your hands! What \emph{(you / do)} \trou{have you been doing}?
  \item He \emph{(write)} \trou{has written} four letters this morning.
  \item It \emph{(rain)} \trou{has been raining} since yesterday and the marshes are under
    water.
  \item Somebody \emph{(eat)} \trou{has eaten} my bread.
  \item I \emph{(seldom see)} \trou{have seldom seen} such a fog on the river.
  \item The police \emph{(fail)} \trou{have failed} to catch the man.
  \item She \emph{(know)} \trou{has known} him since 1880.
  \item How many times \emph{(you / see)} \trou{have you seen} him this week?
\end{enumerate}
\end{anglais}

\note{Forme en \emph{-ing} : 2, 3, 5. Forme simple : 1, 4, 6, 7, 8, 9, 10.}

\note[Trois obstacles]{Phrase 9 : \emph{know} est un verbe d'état, donc pas de \emph{-ing},
même au present perfect. Phrases 1, 4, 10 : dès qu'on compte (\emph{three times}, \emph{four
letters}, \emph{how many times}), la forme simple s'impose, parce qu'on regarde le résultat
et non l'activité. Phrases 2 et 3 : le résultat visible (yeux fatigués, mains sales) sert
justement de preuve d'une activité prolongée.}

%% ===================================================================
\exercice[12 min]{Le paysage et la peur}
\consigne{a) Associez chaque mot à sa traduction.}

\begin{multicols}{2}
\begin{enumerate}[label=\arabic*.,itemsep=1.2mm,leftmargin=*]
  \item a marsh \hfill \trou{e}
  \item a churchyard \hfill \trou{j}
  \item a mound \hfill \trou{b}
  \item cattle \hfill \trou{h}
  \item a gibbet \hfill \trou{a}
  \item a beacon \hfill \trou{l}
  \item to shiver \hfill \trou{d}
  \item to limp \hfill \trou{i}
  \item to growl \hfill \trou{c}
  \item to glare \hfill \trou{k}
  \item to seize \hfill \trou{g}
  \item numb \hfill \trou{f}
\end{enumerate}
\columnbreak
\begin{enumerate}[label=\alph*.,itemsep=1.2mm,leftmargin=*]
  \item un gibet
  \item un tertre
  \item gronder
  \item frissonner
  \item un marais
  \item engourdi
  \item saisir, empoigner
  \item le bétail
  \item boiter
  \item un cimetière
  \item lancer des regards furieux
  \item une balise, un phare
\end{enumerate}
\end{multicols}

\note[Deux remarques]{\emph{Cattle} est un pluriel sans singulier : \emph{the cattle are
feeding}. \emph{A churchyard} est le cimetière attenant à une église, à distinguer de
\emph{a cemetery}, cimetière municipal.}

\consigne{b) Complétez avec un mot de la liste.}
\begin{anglais}
\begin{enumerate}[label=\arabic*.,itemsep=2mm,leftmargin=*]
  \item The man \trou{seized} the boy by the chin and \trou{glared} at him.
  \item He was so cold that his legs were \trou{numb}.
  \item The child stood \trou{shivering} between the tombstones.
  \item Wounded in the leg, the convict \trou{limped} towards the wall.
  \item Beyond the \trou{marshes} lay the leaden line of the river.
\end{enumerate}
\end{anglais}

%% ===================================================================
\newpage
\section{Texte}

\noindent\small\textit{Pip, orphelin de sept ans, est élevé par sa sœur dans les marais de
l'estuaire de la Tamise. Il raconte, bien des années plus tard.}\normalsize

\begin{texte}
My father's family name being Pirrip, and my Christian name Philip, my infant tongue could
make of both names nothing longer or more explicit than Pip. So, I called myself Pip, and
came to be called Pip.

I give Pirrip as my father's family name, on the authority of his tombstone and my sister
--- Mrs. Joe Gargery, who married the blacksmith. As I never saw my father or my mother, and
never saw any \gl[a likeness]{likeness}{un portrait, une image} of either of them (for their
days were long before the days of photographs), my first fancies regarding what they were
like were unreasonably derived from their tombstones. The shape of the letters on my
father's, gave me an odd idea that he was a square, stout, dark man, with curly black hair.
From the character and turn of the inscription, ``\emph{Also Georgiana Wife of the Above},''
I drew a childish conclusion that my mother was \gl[freckled]{freckled}{couvert de taches de
rousseur} and sickly. To five little stone \gl[a lozenge]{lozenges}{un losange}, each about a
foot and a half long, which were arranged in a neat row beside their grave, and were sacred
to the memory of five little brothers of mine --- who gave up trying to get a living,
exceedingly early in that universal struggle --- I am indebted for a belief I religiously
entertained that they had all been born on their backs with their hands in their
trousers-pockets, and had never taken them out in this state of existence.

Ours was the marsh country, down by the river, within, as the river wound, twenty miles of
the sea. My first most vivid and broad impression of the identity of things seems to me to
have been gained on a memorable \gl[raw]{raw}{cru, glacial et humide} afternoon towards
evening. At such a time I found out for certain that this \gl[bleak]{bleak}{désolé, morne}
place overgrown with \gl[a nettle]{nettles}{une ortie} was the churchyard; and that Philip
Pirrip, late of this parish, and also Georgiana wife of the above, were dead and buried; and
that Alexander, Bartholomew, Abraham, Tobias, and Roger, infant children of the aforesaid,
were also dead and buried; and that the dark flat wilderness beyond the churchyard,
intersected with \gl[a dike]{dikes}{un fossé, une digue} and mounds and gates, with scattered
cattle feeding on it, was the marshes; and that the low \gl[leaden]{leaden}{de plomb} line
beyond was the river; and that the distant savage \gl[a lair]{lair}{une tanière} from which
the wind was rushing was the sea; and that the small bundle of shivers growing afraid of it
all and beginning to cry, was Pip.

``Hold your noise!'' cried a terrible voice, as a man started up from among the graves at the
side of the church porch. ``Keep still, you little devil, or I'll cut your throat!''

A fearful man, all in \gl[coarse]{coarse}{grossier, rêche} grey, with a great iron on his
leg. A man with no hat, and with broken shoes, and with an old \gl[a rag]{rag}{un chiffon}
tied round his head. A man who had been soaked in water, and
\gl[to smother]{smothered}{étouffer ; ici : couvrir entièrement} in mud, and
\gl[to lame]{lamed}{estropier} by stones, and cut by \gl[a flint]{flints}{un silex}, and
stung by nettles, and torn by \gl[a briar]{briars}{une ronce}; who limped, and shivered, and
glared, and growled; and whose teeth chattered in his head as he seized me by the chin.

``Oh! Don't cut my throat, sir,'' I pleaded in terror. ``Pray don't do it, sir.''

``Tell us your name!'' said the man. ``Quick!''

``Pip, sir.''

[\ldots]

The man, after looking at me for a moment, turned me upside down, and emptied my pockets.
There was nothing in them but a piece of bread. When the church came to itself --- for he was
so sudden and strong that he made it go head over heels before me, and I saw the steeple under
my feet --- when the church came to itself, I say, I was seated on a high tombstone,
trembling while he ate the bread \gl[ravenously]{ravenously}{voracement}.
\end{texte}
\source{Charles Dickens, \emph{Great Expectations}, Londres, Chapman \& Hall, 1861,
ch.~\textsc{i}.}

\partiel[15 min]{Questions}
\consigne{\en{Answer in English, in full sentences.}}

\begin{enumerate}[label=\arabic*.,itemsep=2mm,leftmargin=*]
  \item \en{Where does the name Pip come from?}
    \lignes{2}
    \corr{From the narrator himself: as a small child he could not pronounce Philip Pirrip
    and made Pip of it, so he called himself Pip and everyone else followed.}
  \item \en{How did Pip picture his parents, and why?}
    \lignes{3}
    \corr{He never saw them and there were no photographs, so he built them out of the
    lettering on their tombstones: a square, stout, dark man with curly hair, and a freckled,
    sickly woman.}
  \item \en{What happened to Pip's five brothers, and how does the narrator speak of them?}
    \lignes{3}
    \corr{They died as infants. The grown-up narrator speaks of them with black humour: they
    \emph{gave up trying to get a living} very early, and he imagines them born with their
    hands in their pockets.}
  \item \en{The long sentence beginning ``At such a time I found out for certain'' names one
    thing after another. What is the last one, and what is the effect?}
    \lignes{3}
    \corr{The last one is Pip himself, reduced to \emph{the small bundle of shivers growing
    afraid of it all and beginning to cry}. The child discovers the world and his own place
    in it in the same movement, and he turns out to be the smallest thing in the landscape.}
  \item \en{List four things we learn about the man from his appearance alone.}
    \lignes{3}
    \corr{The iron on his leg shows he is an escaped convict; he has no hat and broken shoes,
    so he is destitute; he has been soaked, covered in mud, cut and stung, so he has crossed
    the marshes; he is starving, since he eats the boy's bread ravenously.}
  \item \en{Find in the text three words or phrases that make the landscape hostile.}
    \lignes{2}
    \corr{\emph{bleak place overgrown with nettles}, \emph{dark flat wilderness}, \emph{low
    leaden line}, \emph{distant savage lair}, \emph{the wind was rushing}.}
\end{enumerate}

\note[Pluperfect dans le texte]{\emph{A man who had been soaked in water, and smothered in
mud, and lamed by stones} : le pluperfect (au passif) remonte à ce qui a précédé la
rencontre. Plus loin, \emph{they had all been born\ldots and had never taken them out}
remonte à la naissance des frères. Le récit, lui, reste au prétérit.}

%% ===================================================================
\newpage
\section{Version}
\consigne{Traduisez le passage en français. Il vient à la fin du même chapitre : le forçat a
exigé de Pip une lime et de la nourriture, puis s'éloigne.}

\begin{version}
When he came to the low church wall, he got over it, like a man whose legs were numbed and
stiff, and then turned round to look for me. When I saw him turning, I set my face towards
home, and made the best use of my legs. But presently I looked over my shoulder, and saw him
going on again towards the river, still hugging himself in both arms, and picking his way
with his sore feet among the great stones dropped into the marshes here and there, for
stepping-places when the rains were heavy or the tide was in.

The marshes were just a long black horizontal line then, as I stopped to look after him; and
the river was just another horizontal line, not nearly so broad nor yet so black; and the sky
was just a row of long angry red lines and dense black lines intermixed. On the edge of the
river I could faintly make out the only two black things in all the
\gl[a prospect]{prospect}{une vue, un paysage} that seemed to be standing upright; one of
these was the beacon by which the sailors steered --- like an
\gl[unhooped]{unhooped}{sans cercles} cask upon a pole --- an ugly thing when you were near
it; the other, a gibbet, with some chains hanging to it which had once held a pirate. The man
was limping on towards this latter, as if he were the pirate come to life, and come down, and
going back to hook himself up again.
\end{version}
\source{\emph{Ibid.}}

\lignes{22}

\corr{\textbf{Proposition de traduction.} Parvenu au petit mur du cimetière, il l'enjamba
comme un homme dont les jambes seraient engourdies et raides, puis se retourna pour me
chercher des yeux. Le voyant se retourner, je mis le cap sur la maison et fis le meilleur
usage de mes jambes. Mais au bout d'un instant je regardai par-dessus mon épaule et le vis
qui repartait vers la rivière, se serrant toujours dans ses deux bras et cherchant son chemin
de ses pieds meurtris parmi les grosses pierres jetées çà et là dans le marais, pour servir
de gué quand les pluies étaient fortes ou que la marée montait.

Les marais n'étaient plus alors qu'une longue ligne noire horizontale, tandis que je
m'arrêtais à le suivre des yeux ; la rivière n'était qu'une autre ligne horizontale, loin
d'être aussi large ni aussi noire ; et le ciel n'était qu'une rangée de longues raies rouges
et furieuses mêlées de raies noires et épaisses. Au bord de la rivière, je distinguais
faiblement les deux seules choses noires, dans tout ce paysage, qui parussent se tenir
debout : l'une était la balise sur laquelle se dirigeaient les marins --- un tonneau sans
cercles au bout d'une perche, chose bien laide quand on en approchait ; l'autre, un gibet,
auquel pendaient des chaînes qui avaient jadis retenu un pirate. L'homme s'éloignait en
boitant vers ce dernier, comme s'il eût été le pirate revenu à la vie, descendu du gibet, et
s'en retournant s'y accrocher.}

\note[Choix de traduction 1]{\emph{like a man whose legs were numbed and stiff} : la
relative en \emph{whose} se rend par \og dont \fg, mais l'ordre change --- l'anglais place
le nom possédé juste après \emph{whose}, le français le fait suivre du verbe (\og dont les
jambes étaient\ldots \fg).}

\note[Choix de traduction 2]{\emph{When I saw him turning} : \emph{see} + complément +
\emph{-ing} donne en français soit un participe (\og le voyant se retourner \fg), soit une
relative (\og je le vis qui se retournait \fg). Même tour trois lignes plus bas :
\emph{saw him going on again} $\rightarrow$ \og le vis qui repartait \fg. Éviter \og je le
vis se retournant \fg, qui n'est pas français.}

\note[Choix de traduction 3]{\emph{not nearly so broad nor yet so black} : \emph{not nearly}
= \og loin de \fg, et non \og pas près \fg. \emph{Nor yet} renforce la négation : \og ni
même \fg, \og ni non plus \fg.}

\note[Choix de traduction 4]{\emph{as if he were the pirate come to life} : \emph{as if} +
prétérit modal (\emph{were} et non \emph{was}) marque l'irréel ; le français demande le
subjonctif (\og comme s'il eût été \fg{} ou, plus simplement, \og comme s'il avait été \fg).
\emph{Come to life} est un participe passé apposé, et non un impératif.}

\note[Choix de traduction 5]{\emph{an ugly thing when you were near it} : \emph{you} est ici
le \og on \fg{} indéfini du français, comme \emph{one} dans la version de la fiche 1. Le
traduire par \og tu \fg{} ou \og vous \fg{} serait un contresens de registre.}

%% ===================================================================
\partiel{Lexique à mémoriser}
\begin{itemize}[label=--,itemsep=0.8mm,leftmargin=*]\small
  \item \textbf{a marsh} — un marais ; \textit{adjectif} marshy
  \item \textbf{a churchyard} — un cimetière (autour d'une église)
  \item \textbf{a tombstone} — une pierre tombale ; tomb + stone\textit{, le} b \textit{ne se prononce pas}
  \item \textbf{a grave} — une tombe ; \textit{d'où} a graveyard
  \item \textbf{the tide} — la marée ; \textit{parent de l'allemand} Zeit\textit{, le sens premier est « le temps »}
  \item \textbf{to shiver} — frissonner
  \item \textbf{to tremble} — trembler
  \item \textbf{to limp} — boiter
  \item \textbf{to growl} — gronder, grogner ; \textit{formation imitative}
  \item \textbf{to glare at} — fusiller du regard ; \textit{même racine que} glass\textit{, l'idée d'éclat}
  \item \textbf{to seize} — saisir, empoigner
  \item \textbf{to plead} — supplier ; \textit{d'où} a plea \textit{« une supplique »}
  \item \textbf{stout} — corpulent, robuste
  \item \textbf{sore} — douloureux, meurtri
  \item \textbf{numb} — engourdi ; \textit{le} b \textit{est muet, comme dans} thumb, climb
  \item \textbf{stiff} — raide
  \item \textbf{to hug} — serrer dans ses bras
  \item \textbf{to make out} — distinguer, déchiffrer
  \item \textbf{faintly} — faiblement ; faint \textit{« faible, défaillant », d'où} to faint \textit{« s'évanouir »}
  \item \textbf{upright} — droit, debout
\end{itemize}

\end{document}
